\section{PSB işlem hattı: yol, bölme, dengeleme (Path, Split, Balance)}\label{sec:method}

PSB; önce rota, sonra küme yaklaşımına dayanan bir kurulum adımından \citep{beasley1983,prins2004},
bunu izleyen ve min--maks amacına özgü bir yerel aramadan ve her an sonuç veren bir geniş komşuluk
aramasından (LNS) oluşmaktadır. Tüm hesaplama çekirdekleri Numba \citep{lam2015numba} ile derlenmiştir. Rota başına
çalışan hesaplama çekirdekleri yorumlayıcı kilidini serbest bıraksa da bunların paralel çalıştırılması,
\nameref{sec:results} bölümünde raporlanan ölçümlerde kazanç sağlamamıştır; bu nedenle PSB, LKH-3 ve diğer
çözücülerde olduğu gibi tüm deneylerde \emph{tek çekirdek} üzerinde çalıştırılmıştır. İşlem hattı Algoritma~\ref{alg:psb}'de
özetlenmiştir.

\begin{algorithm}[t]
\caption{Her an sonuç veren LNS ile PSB (Path--Split--Balance)}\label{alg:psb}
\begin{algorithmic}[1]
\Require hedefler $P$, depo $D$, robot sayısı $k$, zaman bütçesi $B$
\State $N \gets$ $K$-en yakın komşu aday listeleri ($K=10$, k-d ağacı)
\State $\pi \gets \textsc{GreedyEdge}(P,N)$; \; $\pi\gets\textsc{TwoOptOrOpt}(\pi,N)$ \Comment{dev tur (Path)}
\State $\{R_j\} \gets \textsc{OptimalSplit}(\pi, D, k)$ \Comment{bölme (Split), Lemma~\ref{lem:split}}
\State tüm $j$ için $R_j \gets \textsc{TwoOptOrOpt}(\{D\}\cup R_j)$ \Comment{cilalama, her rota için bağımsız}
\Repeat \Comment{dengeleme (Balance)}
  \State $\max_j L(R_j)$ değerini düşüren en iyi taşıma/takas hamlesini uygula; etkilenen rotaları yeniden cilala
\Until{iyileştiren hamle kalmayana dek}
\While{geçen süre $< B$} \Comment{her an sonuç veren LNS}
  \State $j^\star\gets\arg\max_j L(R_j)$; ağırlık merkezi bu rotaya en yakın $m-1$ rotayı seç
  \State $m$-robotlu alt problemi 2--7. satırlarla yeniden çöz (rastgele başlangıç)
  \State alt problemin tamamlanma süresi azalırsa kabul et; durağanlıkta $m$'yi büyüt
\EndWhile
\end{algorithmic}
\end{algorithm}

\subsection{Path: Greedy-Edge ile dev tur}
Dev tur, Greedy-Edge (GE) sezgiseli \citep{bentley1992} ile oluşturulmuştur:
$K$-en yakın komşu listelerinden gelen aday kenarlar artan uzunluk sırasıyla taranmış;
bir kenar, her iki uç noktasının derecesi $<2$ olduğunda ve bir çevrimi kapatmadığında
(union--find) kabul edilmiştir. Ortaya çıkan yol parçaları en yakın serbest uç noktalar
üzerinden açgözlü biçimde birleştirilmiştir. GE dönmeye karşı değişmezdir ve kafes benzeri
örnekler için yapıcı sezgiseller üzerine yürüttüğümüz tamamlayıcı çalışmada (hakem değerlendirmesinde) en
dayanıklı tek tur yapıcı yöntem olarak öne çıkmıştır. Tur daha sonra aday listeleriyle sınırlandırılmış 2-opt \citep{croes1958} ve Or-opt \citep{or1976}
(uzunluğu 1--3 olan segmentler, her iki yönelimde) ile, konum indeksli bir dizi gösterimi
kullanılarak iyileştirilmiştir. GE'nin maliyeti $O(nK\log(nK))$ olup uygulamada
sabit sayıda yerel arama geçişi kullanılmıştır; yerel arama ve dengeleme yineleme
sayılarına veri bağımsız bir üst sınır konmadığından bu, tüm işlem hattı için
kötü durum karmaşıklığı iddiası değildir.

\subsection{Split: dev turun optimal min--maks bölünmesi}
$\pi=(p_0,\dots,p_{n-1})$ döngüsel olarak okunan dev tur olsun. Ardışık $p_s,\dots,p_e$
bloğundan oluşan bir rotanın uzunluğu
$c(s,e)=d(D,p_s)+\sum_{t=s}^{e-1}d(p_t,p_{t+1})+d(p_e,D)$ olarak verilir.

\begin{lemma}\label{lem:split}
Sabit bir döngüsel sıra için (i) $c(s,e)$, $e$'ye göre azalmayan ve $s$'ye göre artmayan
bir fonksiyondur; dolayısıyla (ii) bir $T$ eşiği için her rotayı $c\le T$ koşulunun izin verdiği
ölçüde açgözlü biçimde uzatmak, herhangi bir $s$ başlangıcından itibaren gereken rota sayısını en aza indirir ve
(iii) en az bir başlangıç için en fazla $k$ rotanın yettiği en küçük $T$ değeri, $\pi$'nin
optimal min--maks bölmesini verir.
\end{lemma}
\begin{proof}
(i) Üçgen eşitsizliği gereği $c(s,e+1)-c(s,e)=d(p_e,p_{e+1})+d(p_{e+1},D)-d(p_e,D)\ge 0$
olur; $s$ için argüman simetriktir. (ii) (i) gereği, olurlu bir
çözüm ilk rotasını $e'$ noktasında bitiriyorsa açgözlü bitiş noktası için $e\ge e'$ olur ve açgözlü çözümün kalan
son eki diğer çözümün son ekinin içinde kalır; rota sayısı üzerinden tümevarımla
iddia elde edilir. (iii), olurluluğun $T$'ye göre monoton olmasından
çıkar.
\end{proof}

Önek toplamları ve iki katına çıkarılmış dizi üzerinde iki işaretçili bir taramayla her başlangıcın
açgözlü ardılı $\mathrm{nxt}(s)$ $O(n)$ sürede hesaplanmıştır; $n$ başlangıcın tamamının
$k$ atlamayla sınanması $O(nk)$ maliyetindedir ve $T$ üzerinde $10^{-7}$ göreli duyarlığa kadar
yapılan ikiye bölme araması yaklaşık 25 sınama gerektirir. Dolayısıyla bölme, dev tur \emph{verildiğinde}
optimaldir ve $n=10^4$ için bile milisaniyeler içinde hesaplanır.

\subsection{Cilalama ve dengeleme}
Depo eklendiğinde her rota küçük ve kapalı bir TSP turuna dönüşür; bu turlar aynı 2-opt/Or-opt hesaplama çekirdeğiyle
birbirinden bağımsız olarak cilalanmıştır (doğrudan paralelleştirilebilir; yeterli çekirdek bulunduğunda her rota için
bir iş parçacığı). Cilalama hiçbir rotanın uzunluğunu artırmaz;
ancak bölme sınırları cilalanmamış uzunluklara göre seçildiğinden rotalar
hafifçe dengesizleşir. \textsc{Balance} adımı bu dengesizliği, yalnızca o anda en uzun olan $R_{j^\star}$ rotası için
ve yalnızca bu rotanın her köşesinin $K$ en yakın komşusuna yönelik değerlendirilen iki rotalar arası
hamleyle giderir: \emph{taşıma} (relocate; $v\in R_{j^\star}$ düğümünü
$R_{j'}$ rotasındaki bir $c$ komşusunun yanına, $c$'den önce ya da sonra taşıma) ve \emph{takas}
(swap; $v$ ile $c$'nin yer değiştirmesi). Bir hamle,
$\max\{L(R_{j^\star})-\Delta^-,\;L(R_{j'})+\Delta^+\}<L(R_{j^\star})$ ise, yani etkilenen
çiftin tamamlanma süresi kesin olarak azalıyorsa uygulanmıştır. Rotalar, her robot için bir depo nöbetçi
düğümü içeren çift bağlı bir listede saklanmış, böylece her hamlenin maliyeti $O(1)$ olmuştur; değişen
rotalar yeniden cilalanmış ve dengeleme--cilalama döngüsü iyileşme kalmayana dek
sürdürülmüştür.

\subsection{Her an sonuç veren LNS}
Yukarıda açıklanan biçimiyle PSB, tek atış bir yapıcı yöntemdir. Ek süreden yararlanmak için LNS,
yinelemeli olarak en uzun rotayı ve ağırlık merkezi bu rotaya en yakın olan $m-1$ rotayı
($m\in\{2,3,4\}$, rastgele) seçmiş, bunların hedeflerini birleştirmiş ve ortaya çıkan
$m$-robotlu alt problemi PSB işlem hattının tamamıyla yeniden çözmüştür. Çeşitlendirme,
alt problemin dev turu aracılığıyla sağlanmıştır: bu tur ya GE ile yeniden oluşturulmuş ya da
mevcut rotaların depo çevresinde rastgele bozulmuş bir açısal sırayla art arda eklenmesiyle
(her rota rastgele ters çevrilerek) başlatılmış ve ardından 2-opt/Or-opt ile iyileştirilmiştir. Yeni
rotalar eskilerinin yerini yalnızca alt problemin tamamlanma süresi kesin olarak azaldığında almıştır;
böylece genel tamamlanma süresi monoton (artmayan) kalmıştır. Art arda $50k$ başarısız denemeden sonra
komşuluk genişletilmiş ($m\gets m+1$); komşuluk $k$ rotanın tamamını kapsadıktan sonra
yeniden $50k$ kez başarısız olunduğunda LNS durdurulmuştur. Kabul edilen her iyileşme
zaman damgasıyla kaydedilmiş, böylece aşağıda raporlanan, süreye bağlı iyileşme eğrileri elde edilmiştir.

\subsection{Hasat işlemleri için uzantılar}\label{sec:ext}
Hasat uygulaması üç değişiklik gerektirmiş;
Path--Split--Balance yapısı ise korunmuştur. Kapasiteli çok seferler ve olay
tetiklemeli yeniden planlamayı içeren bu varyant \emph{PSB-MT/R}
(çok seferli/yeniden planlamalı PSB) olarak adlandırılmıştır.

\paragraph{Kapasiteli seferler} Bir hasat platformu tek bir kasa taşır ve kasa her dolduğunda
kasa istasyonuna döner; dolayısıyla her platform bir dizi
\emph{sefer} gerçekleştirir. Dev tur verildiğinde, tahmini yükü kasa kapasitesi $Q$'yu
aşmayan ve toplam süresi (yol, hizmet ve sefer başına bir kasa değişimi) en küçük olan
seferlere bölme işlemi, dev tur sırası üzerinde bir en kısa yol problemine karşılık gelir; bu problem, $L$ sefer başına
en fazla durak sayısı olmak üzere, dinamik programlamayla $O(nL)$ sürede kesin olarak
çözülmüştür \citep{beasley1983,prins2004}. Seferler daha sonra
her platformun hazır olma zamanını gözeten en uzun işlem süresi kuralıyla $k$ platforma atanmış,
ardından en yüklü platform ile diğerleri arasında taşıma ve takas adımları
(dengeleme) uygulanmıştır.

\paragraph{Bahçe eksenine duyarlı dev tur} Platformlar yalnızca ağaçlar arasındaki sıra aralarında
(koridorlarda) hareket edebildiğinden yol uzunluğu, bahçenin kendi eksen takımında ölçülmüştür: Manhattan
($L_1$) uzaklığına, her iki koordinat birlikte değiştiğinde bir dönüş cezası eklenmiştir. Bahçe
ekseni planlayıcıya verilmemiş, ağaç konumlarından kestirilmiştir.
Kaba açı, her ağaçtan en yakın dört komşusuna uzanan yönlerin $90^\circ$ periyotla
alınan dairesel ortalaması olarak hesaplanmıştır (kare ya da
dikdörtgen bir kafeste en yakın komşular eksenler boyunca yer alır, köşegen komşular ise
$\sqrt2$ kat daha uzaktadır; bu nedenle kestirim köşegen doğrultulara kilitlenmez). Açı
daha sonra, izdüşürülen konumların tarak skoru en büyüklenerek $\pm5^\circ$ aralığında inceltilmiştir.
Tarak skoru, her eksendeki izdüşüm histogramının normalize edilmiş kütleleri
$h_b$ için $B\sum_b h_b^2$ ($B$: kutucuk sayısı) olup sıra izdüşümleri birkaç
dar seviyede toplandığında büyür \citep{asil2026frame}. Bu eksen takımında dev tur, ağaç dizilerinin yılan biçimli
(boustrophedon) sıralamasıdır; her iki dizi yönü denenmiş ve
planlanan tamamlanma süresi daha küçük olan seçilmiştir. Sıralı (telli terbiye) bahçelerde aynı
kurulum sıra araları boyunca uygulanmış; bir sıra arasından diğerine geçişte ise yol, sıra arası uçlarının
dışbükey zarfı çevresindeki baş tarlayı izlemiştir.

\paragraph{Olay tetiklemeli yeniden planlama} Bir platform kasa değişimi için her döndüğünde,
kalan iş, ağaç başına olgun meyveye ilişkin en güncel kestirimler kullanılarak tüm filo için
yeniden planlanmıştır (bkz.\ hasat benzetimi); hâlihazırda sürmekte olan
seferler değiştirilmemiştir. Planlama, kasa değiştirilirken tek çekirdekte yürütülmüştür; dolayısıyla
bir plan, platformu yalnızca hesaplanması kasa değişiminden uzun sürdüğünde geciktirir.

\paragraph{Uç devralmalı k-d bölgeleri} Merkezi olmayan bir alternatif olarak tahmini
iş yükü, kestirilen eksenler doğrultusunda bir $k$-d bölüntüsüyle bir kez $k$ adet derli toplu bölgeye ayrılmış,
kesme çizgileri sıra aralarına yerleştirilmiştir; her platform
kendi bölgesini yılan biçimli sırayla işlemiştir. Bir platform kendi bölgesini bitirdiğinde,
kalan işi en fazla olan platformun geri kalan dizisini devralır: dizi, iki platformun
aynı anda bitirmesinin beklendiği noktadan kesilir ve yardımcı platform,
kesilen kısmın uzak ucundan başlayarak bölgenin asıl sahibi olan platforma doğru ilerler. Bu yaklaşım, bahçe geometrisine uyarlanmış
bir iş çalma \citep{blumofe1999} biçimidir; merkezi yeniden planlama gerektirmez ve platformları
karşılaşıncaya kadar tarlanın ayrık bölümlerinde tutar.

\paragraph{Yapay zekâ araçlarının kullanımı} Yazar, metnin ve bilimsel anlatımın
gözden geçirilmesinde yapay zekâ destekli bir araçtan düzenleme önerileri almak
amacıyla yararlanmıştır. Yöntemin ve deneylerin tasarımı, veri analizi ve bilimsel
sonuçların doğrulanması yazara aittir; nihai metin yazar tarafından incelenmiş ve
onaylanmıştır.
